%Optimizability: given a transformer, does there exist a (descriptionally) smaller one 
%
%\fix{Maybe we need to define it as a function problem and then show that equivalence is hard.}

As a consequence of our results, we can show that reasoning about languages of UHATs (e.g.
equivalence, emptiness, universality, etc.) is provably intractable. 
Contrast this to deterministic finite automata, where all these problems can be done in
polynomial time (cf. \citep{Kozen}). In the following,
%Using our previous results, 
we show the precise complexity for the \emph{equivalence problem}, 
i.e., the problem of checking whether two given UHATs recognize the same 
language.
That the universality problem is $\EXPSPACE$-complete can be proven similarly.

\begin{theorem}
Equivalence of UHATs is $\EXPSPACE$-complete.
\end{theorem}
\begin{proof}
To prove the lower bound, we reduce from the non-emptiness problem for UHATs, 
which by \cref{thm:UHAT-emptiness} is $\EXPSPACE$-complete.
To this end, let $\mathcal{T}$ be a given UHAT and fix a UHAT $\mathcal{T}_0$ that recognizes the empty language. 
Then we have that $\mathcal{T}$ and $\mathcal{T}_0$ are equivalent if and only if $\mathcal{T}$ recognizes the empty language.

For the upper bound let the UHATs $\mathcal{T}_1$ and $\mathcal{T}_2$ be given.
We apply \cref{prop:UHAT-LTL} to turn $\mathcal{T}_1$ and $\mathcal{T}_2$ in exponential time
into LTL formulas $\varphi_1$ and $\varphi_2$, respectively.
Now, $\mathcal{T}_1$ and $\mathcal{T}_2$ are equivalent if and only if $\varphi_1$ and $\varphi_2$ are equivalent.
The latter can be decided in polynomial space \citep{SistlaC85}, which results in an exponential-space algorithm in total.
\end{proof}